\documentclass[11pt,a4paper]{article}

% ============================================================
% PAQUETES
% ============================================================
\usepackage[utf8]{inputenc}
\usepackage[T1]{fontenc}
\usepackage[spanish,es-noshorthands]{babel}
\usepackage{amsmath,amssymb}
\usepackage{graphicx}
\usepackage{geometry}
\geometry{top=2.2cm, bottom=2.2cm, left=2.2cm, right=2.2cm}
\usepackage{float}
\usepackage{enumitem}
\setlist{nosep, leftmargin=1.5em}
\usepackage{parskip}
\setlength{\parindent}{1.25em} % Sangría de primera línea
\usepackage{indentfirst} % Sangra también el primer párrafo tras cada sección
\usepackage{caption}
\usepackage{url}
\usepackage{xcolor}
\usepackage{booktabs}
\usepackage{tabularx}
\usepackage{array}
\newcolumntype{L}{>{\raggedright\arraybackslash}X}

% Configuración de hyperref (combinando el color azul de enlaces con el estilo sin marcos)
\usepackage{hyperref}
\hypersetup{
	colorlinks=true,
	linkcolor=black,
	urlcolor=blue,
	citecolor=black,
	pdfborder={0 0 0}
}

% Renombrar cuadros a Tablas en español
\addto\captionsspanish{\renewcommand{\tablename}{Tabla}}

% Títulos de sección más compactos
\makeatletter
\renewcommand\section{\@startsection{section}{1}{\z@}%
  {-3ex \@plus -.5ex \@minus -.2ex}{1.2ex \@plus .2ex}{\normalfont\large\bfseries}}
\renewcommand\subsection{\@startsection{subsection}{2}{\z@}%
  {-2.2ex \@plus -.5ex \@minus -.2ex}{0.8ex \@plus .2ex}{\normalfont\normalsize\bfseries}}
\makeatother

% Entorno thebibliography sin el encabezado automático "Referencias":
% ya ponemos nuestro propio \section*{Bibliografía} antes de usarlo.
\makeatletter
\renewenvironment{thebibliography}[1]
 {\list{\@biblabel{\@arabic\c@enumiv}}%
       {\settowidth\labelwidth{\@biblabel{#1}}%
        \leftmargin\labelwidth
        \advance\leftmargin\labelsep
        \itemsep 1.5pt \parsep 0pt
        \usecounter{enumiv}%
        \let\p@enumiv\@empty
        \renewcommand\theenumiv{\@arabic\c@enumiv}}%
  \sloppy\clubpenalty4000\widowpenalty4000%
  \sfcode`\.\@m}
 {\def\@noitemerr
   {\@latex@warning{Empty `thebibliography' environment}}%
  \endlist}
\makeatother
\captionsetup{
    font={small,it},
    labelfont={small,it,bf},
    justification=centering,
    singlelinecheck=false
}

% ============================================================
% DATOS DEL TRABAJO
% ============================================================
\newcommand{\materia}{Procesamiento de Señales e Imágenes Biomédicas}
\newcommand{\codigoMateria}{16.63}
\newcommand{\tituloTrabajo}{Anteproyecto}
\newcommand{\subtituloTrabajo}{Detección de lunares sospechosos mediante imágenes dermatoscópicas}
\newcommand{\integranteUno}{65555 -- Cassano Nieto, Facundo Tomás (\href{mailto:fcassanonieto@itba.edu.ar}{fcassanonieto@itba.edu.ar})}
\newcommand{\integranteDos}{65414 -- D'Elía, Lucas Valentín (\href{mailto:ldelia@itba.edu.ar}{ldelia@itba.edu.ar})}
\newcommand{\integranteTres}{65560 -- Ortega, Facundo (\href{mailto:facortega@itba.edu.ar}{facortega@itba.edu.ar})}
\newcommand{\comision}{BIO}
\newcommand{\docenteUno}{Tomás, Roberto Sebastián}
\newcommand{\docenteDos}{Capurro, Lucila}
\newcommand{\docenteTres}{Reppucci, Victoria}
\newcommand{\docenteCuatro}{Tassara, Francisco Joaquín}
\newcommand{\fechaEntrega}{10 de octubre de 2026}
\newcommand{\logoITBA}{logo_itba.png} % colocar el archivo del logo en esta carpeta con este nombre

\begin{document}

% ============================================================
% CARÁTULA
% ============================================================
\begin{titlepage}
    \centering
    \vspace*{1cm}
    \IfFileExists{\logoITBA}{%
        \includegraphics[width=0.45\textwidth]{\logoITBA}\par
        \vspace{0.1cm}
    }{}
    {\LARGE Instituto Tecnológico de Buenos Aires\par}
    \vspace{0.3cm}
    {\Large \materia\ (\codigoMateria)\par}
    \vspace{0.3cm}
    {\large Comisión \comision\par}
    \vfill
    {\Huge\bfseries \tituloTrabajo\par}
    \vspace{0.15cm}
    {\LARGE \subtituloTrabajo\par}
    \vfill
    {\large\bfseries Grupo 2\par}
    \vspace{0.3cm}
    {\large \integranteUno \par}
    {\large \integranteDos \par}
    {\large \integranteTres \par}
    \vfill
    {\large\bfseries Docentes\par}
    \vspace{0.3cm}
    {\large \docenteUno \par}
    {\large \docenteDos \par}
    {\large \docenteTres \par}
    {\large \docenteCuatro \par}
    \vfill
    {\large \fechaEntrega \par}
\end{titlepage}


% ============================================================
% 1. INTRODUCCIÓN
% ============================================================
\section{Introducción}

\subsection{Problemática}

El melanoma es un cáncer de piel, y su pronóstico depende de cuán
temprano se lo detecte. Al comienzo, se parece mucho a un lunar común, y
el dermatólogo lo evalúa con el dermatoscopio, que posee una lente de
aumento con luz propia. Para hacerlo de forma sistemática usa la regla
ABCD, que puntúa la asimetría ($A$, de 0 a 2 según en cuántos ejes es
asimétrica la lesión), el borde ($B$, de 0 a 8 según en cuántos sectores
el pigmento termina de forma abrupta), los colores ($C$, de 1 a 6) y las
estructuras diferenciales, como la red de pigmento o los glóbulos ($D$,
de 1 a 5), y los combina en el puntaje total
\[
\mathrm{TDS} = 1{,}3\,A + 0{,}1\,B + 0{,}5\,C + 0{,}5\,D.
\]
Un TDS menor a 4,75 indica una lesión benigna; entre 4,75 y 5,45, una
lesión sospechosa, y mayor a 5,45, una lesión probablemente maligna
\cite{nachbar}. El único problema es que la evaluación es visual y su
precisión depende de la experiencia de quien la hace.

\subsection{Marco teórico}

\textbf{Imagen dermatoscópica y color.} Es una imagen RGB en la que el
color tiene valor diagnóstico, porque depende de la profundidad del
pigmento: los marrones y el negro indican melanina en la epidermis; el
azul grisáceo, melanina en la dermis, y el blanco, zonas de regresión.
Nos convendría trabajar con el espacio HSV, que separa matiz ($H$) y
saturación ($S$) con respecto al brillo ($V$).

\textbf{Artefactos.} Hay pelos de la piel que cruzan por arriba de la
lesión (líneas oscuras y finas) que desplazan el borde segmentado y
agregan un color que no pertenece a la lesión. Podrían quitarse con
técnicas de morfología como un cierre (dilatación+erosión) con un
elemento estructurante más ancho que un pelo y más chico que la lesión.
También aparecen burbujas por el gel que usa el experto para aumentar la
imagen, y esquinas oscuras por la lente circular (una viñeta), que más
adelante explicamos cómo se podrían mejorar.

\textbf{Textura.} Las estructuras diferenciales son patrones que se
repiten dentro de la lesión. Podríamos usar la matriz de co-ocurrencia
de niveles de gris (GLCM) para contar con qué frecuencia aparecen juntos
dos niveles de gris a una distancia y dirección dadas, y calcular
atributos como la energía, la entropía, el contraste, la homogeneidad y
la correlación para clasificar las texturas.

\subsection{Estado del arte}

Los sistemas de diagnóstico asistido segmentan la lesión, miden
descriptores de forma, color y textura y los entregan a un clasificador
\cite{celebi}. Las redes neuronales convolucionales ya alcanzan un
desempeño comparable al de los dermatólogos \cite{esteva}, pero no
muestran en qué se basa cada decisión, y la mayoría solo informa el
diagnóstico final, por lo que no permiten comprobar si cada criterio se
mide bien.

\subsection{Pregunta, objetivo e hipótesis}

\textbf{Pregunta.} ¿Pueden medirse automáticamente los criterios de la
regla ABCD y alcanzar el puntaje establecido por los dermatólogos para
distinguir lesiones sospechosas de lunares comunes?

\textbf{Objetivo general.} Desarrollar un algoritmo que segmente la
lesión en una imagen dermatoscópica, calcule los criterios de la regla
ABCD y el TDS y estime si la lesión es sospechosa, con el soporte de una
interfaz gráfica que muestre cada etapa.

\textbf{Hipótesis.} 1) Quitar los pelos con un cierre morfológico mejora
la segmentación de las imágenes con pelos respecto de no quitarlos y de
un filtro de mediana. 2) El TDS automático separa las lesiones
sospechosas de los nevos comunes, y agregar textura mejora esa
separación. 3) El desempeño se mantiene aunque apliquemos una compresión
JPEG.

% ============================================================
% 2. MATERIALES Y MÉTODOS
% ============================================================
\section{Materiales y métodos}

\subsection{Base de datos}

Usaremos la base PH$^2$ \cite{ph2}, del Hospital Pedro Hispano
(Portugal). Contiene 200 imágenes dermatoscópicas RGB de
$768 \times 560$ píxeles, adquiridas con el mismo equipo y en
condiciones uniformes, de 80 lunares comunes, 80 nevos atípicos y 40
melanomas.

Incluye el diagnóstico clínico e histológico, la máscara de la lesión
dibujada a mano y la evaluación de un dermatólogo de cada criterio:
asimetría (0, 1 o 2 ejes), presencia de cada uno de los seis colores y
presencia de red de pigmento, puntos y glóbulos, estrías, áreas de
regresión y velo azul-blanquecino. Eso permite validar el algoritmo paso
a paso: la segmentación contra la máscara, $A$ y $C$ contra el
dermatólogo, la textura contra las estructuras anotadas y la
clasificación contra el diagnóstico. Investigaremos en qué criterios
coincide el algoritmo con el médico, en cuáles falla y por qué, y cuánto
aporta cada uno a la clasificación.

\subsection{Metodología}

\begin{enumerate}[label=\arabic*., leftmargin=1.6em, itemsep=6pt, topsep=4pt]
\item \textbf{Análisis exploratorio.} Analizar los histogramas de cada
  canal (RGB y HSV) dentro de la lesión y en la piel, para elegir el
  canal con más contraste; tomar una fracción de cada imagen cubierta por
  pelos, y ver los valores típicos en HSV de cada uno de los seis
  colores, a partir de las máscaras de color de PH$^2$.

\item \textbf{Mejoramiento.} Las esquinas oscuras del viñeteo se detectan
  con un umbral y etiquetando las regiones conectadas al borde de la
  imagen, y se excluyen. Para los pelos, se aplica a cada canal un cierre
  morfológico con elementos estructurantes en todas las orientaciones.

\item \textbf{Compresión y descompresión.} Las imágenes se comprimen con
  JPEG a distintas calidades y se descomprimen, se comparan y se repite
  el algoritmo sobre las imágenes comprimidas.

\item \textbf{Segmentación.} Se comparan el umbral de Otsu,
  \emph{region-growing} desde el centro de la imagen, k-means y
  Chan-Vese. Luego se aplica una apertura para eliminar islas, un cierre
  para unir los huecos del borde y, con el etiquetado, se conserva la
  región que contiene el centro. Para finalizar, nos quedamos con el
  método que dé un mayor Dice en la validación cruzada con los datos de
  PH$^2$.

\item \textbf{Criterios A, B y C.} \emph{Asimetría}: se calculan el
  centroide y los ejes principales de la máscara, se la refleja sobre
  cada eje y se mide la fracción de área que no se superpone. Un eje
  cuenta como asimétrico si esa fracción supera un umbral, elegido con
  las imágenes de entrenamiento para coincidir lo más posible con la
  evaluación del dermatólogo, y $A$ es la cantidad de ejes asimétricos
  (0, 1 o 2). \emph{Borde}: la lesión se divide en 8 sectores y en cada
  uno se calcula, con el filtro de Sobel sobre $V$, el gradiente medio
  sobre el contorno; los sectores con gradiente alto cuentan como borde
  abrupto. \emph{Color}: cada píxel se asigna al más cercano de los seis
  colores de referencia en HSV, y un color está presente si ocupa más de
  un porcentaje mínimo de la lesión. Además, se puede agregar un análisis
  de circularidad $4\pi A/P^2$ (con $A$ el área y $P$ el perímetro del
  lunar).

\item \textbf{Textura y criterio D.} Se hace un análisis con GLCM sobre
  la capa $V$ dentro de la lesión, promediada en las cuatro direcciones,
  y se analizan sus atributos de energía, entropía, contraste, IDM,
  correlación, \emph{cluster shade} y \emph{cluster prominence}. Como
  PH$^2$ no anota exactamente las cinco estructuras de la regla, usamos
  un $D$ adaptado: la cantidad de estructuras que el dermatólogo marcó
  como presentes (red de pigmento, puntos o glóbulos, estrías, áreas de
  regresión y velo). Para estimarlo en una imagen nueva se usa el
  atributo de textura que más cambia con esa cantidad en las imágenes de
  entrenamiento: sobre él se fijan umbrales que separan los distintos
  valores de $D$, y a cada imagen se le asigna el valor del intervalo en
  el que cae.

\item \textbf{Clasificación por umbral.} Una lesión se clasifica como
  sospechosa si su TDS automático supera un umbral. Se comparan los
  umbrales clínicos de la regla (4,75 para lesión sospechosa y 5,45 para
  probable melanoma) con un umbral elegido con las imágenes de
  entrenamiento como el que mejor separa a ambos grupos, es decir, el que
  maximiza la suma de la sensibilidad y la especificidad. Para la
  hipótesis 2, al TDS se le suma el atributo de textura que más
  diferencia a ambos grupos en el entrenamiento, llevado a la misma
  escala (restándole su media y dividiéndolo por su desvío), y se vuelve
  a elegir el umbral.

\item \textbf{Interfaz gráfica.} Aplicación en Streamlit que muestra la
  máscara de pelos, el contorno, los ejes de simetría, los sectores del
  borde, el mapa de colores, los valores de $A$, $B$, $C$ y $D$, el TDS y
  la clasificación resultante.

\item \textbf{Evaluación.} Sobre las imágenes de prueba, la segmentación
  se compara con la máscara de PH$^2$ mediante el coeficiente de Dice (1
  si las máscaras coinciden y 0 si no se superponen). Los criterios $A$
  y $C$ se comparan con el dermatólogo mediante el $\kappa$ de Cohen, que
  mide cuánto coinciden descontando las coincidencias que se darían por
  azar; valores entre 0,41 y 0,60 se consideran un acuerdo moderado
  \cite{landis}. La clasificación se mide con la sensibilidad (fracción
  de lesiones sospechosas detectadas), la especificidad (fracción de
  lunares comunes reconocidos como tales) y el área bajo la curva ROC
  (AUC), que resume ambas para todos los umbrales posibles: vale 1 si el
  puntaje separa perfectamente a los dos grupos y 0,5 si no es mejor que
  el azar.

  Para las hipótesis: 1) en las imágenes con pelos, se compara el Dice
  promedio sin quitar los pelos, con el filtro de mediana y con el cierre
  morfológico, y se cuenta en cuántas imágenes mejora cada método; 2) se
  comparan la sensibilidad, la especificidad y el AUC del TDS solo y del
  TDS con textura; 3) se repite el algoritmo sobre las imágenes
  comprimidas con JPEG a distintas calidades y se grafica cómo cambian el
  Dice, la sensibilidad y la especificidad a medida que baja la calidad.
\end{enumerate}

% ============================================================
% BIBLIOGRAFÍA
% ============================================================
\section*{Bibliografía}
\footnotesize
\begin{thebibliography}{99}
\bibitem{nachbar} F.~Nachbar et al., ``The ABCD rule of dermatoscopy: high prospective value in the diagnosis of doubtful melanocytic skin lesions,'' \emph{J. Am. Acad. Dermatol.}, 30(4):551--559, 1994.
\bibitem{celebi} M.~E.~Celebi et al., ``A methodological approach to the classification of dermoscopy images,'' \emph{Comput. Med. Imaging Graph.}, 31(6):362--373, 2007.
\bibitem{esteva} A.~Esteva et al., ``Dermatologist-level classification of skin cancer with deep neural networks,'' \emph{Nature}, 542:115--118, 2017.
\bibitem{ph2} T.~Mendonça, P.~M.~Ferreira, J.~S.~Marques, A.~R.~S.~Marcal, J.~Rozeira, ``PH$^2$: a dermoscopic image database for research and benchmarking,'' en \emph{Proc. 35th IEEE EMBC}, pp.~5437--5440, 2013.
\bibitem{wang} Z.~Wang, A.~C.~Bovik, H.~R.~Sheikh, E.~P.~Simoncelli, ``Image quality assessment: from error visibility to structural similarity,'' \emph{IEEE Trans. Image Process.}, 13(4):600--612, 2004.
\bibitem{landis} J.~R.~Landis, G.~G.~Koch, ``The measurement of observer agreement for categorical data,'' \emph{Biometrics}, 33(1):159--174, 1977.
\end{thebibliography}

\end{document}
